% Requer pacotes no preâmbulo do documento principal:
% \usepackage{tikz}
% \usetikzlibrary{shapes.geometric, arrows.meta, positioning}

\begin{tikzpicture}[
    node distance=2cm and 1.5cm,
    block/.style={rectangle, draw, fill=blue!5, text width=3.5cm, text centered, rounded corners, minimum height=1.5cm, font=\sffamily\small},
    arrow/.style={-{Stealth[scale=1.2]}, thick, draw=black!70},
    dashed_arrow/.style={-{Stealth[scale=1.2]}, dashed, thick, draw=red!70!black},
    green_arrow/.style={-{Stealth[scale=1.2]}, thick, draw=green!60!black}
]

% Nós principais
\node (input) [block, fill=gray!10] {Imagem Original\\(Câmera/Sensor)};
\node (enhancement) [block, right=2.5cm of input, fill=blue!10] {Técnicas de Aprimoramento\\(Denoising, Sharpening, etc.)};
\node (output) [block, right=2.5cm of enhancement, fill=gray!10] {Imagem Aprimorada};

% Nós de desfecho
\node (human) [block, above=1.5cm of output, fill=green!10, text width=3cm] {Observador Humano\\(Qualidade Perceptual)};
\node (cnn) [block, below=1.5cm of output, fill=red!10, text width=3cm] {Classificador CNN\\(Desempenho)};

% Nó explicativo
\node (spectrum) [block, below=1.2cm of enhancement, fill=yellow!15, text width=4cm] {Alteração no Espectro\\(Perda de Altas Frequências e Sinais Discriminativos)};

% Conexões principais
\draw [arrow] (input) -- (enhancement);
\draw [arrow] (enhancement) -- (output);

% Conexões para Humano
\draw [green_arrow] (output) -- (human) node[midway, right, font=\sffamily\scriptsize, text=green!50!black] {Efeito Positivo};

% Conexões para CNN
\draw [dashed_arrow] (output) -- (cnn) node[midway, right, font=\sffamily\scriptsize, text=red!50!black] {Degradação Sistemática};

% Conexões da alteração de espectro
\draw [arrow] (enhancement) -- (spectrum) node[midway, right, font=\sffamily\scriptsize] {Causa};
\draw [dashed_arrow] (spectrum) -- (cnn) node[midway, below, font=\sffamily\scriptsize, text=red!50!black, sloped, align=center] {Redução de Features\\Relevantes};

\end{tikzpicture}
